A new family of programmable payment protocols (x402, AP2, A2A--x402, A402) lets autonomous agents pay for online services at machine speed, but supplies no algorithm for the buyer's first decision: which provider to pay. We formalize this gap as \emph{agent-native bandits}---budget-aware, non-stationary, contextual decision-making over heterogeneous paid providers, where the bandit is the autonomous buyer itself rather than a human-operated routing service. We propose \emph{402Pilot}, a protocol-agnostic decision layer with stable interfaces between agent decision-making and payment execution, and \emph{PA-DCT} (Payment-Aware Discounted Contextual Thompson sampling), the policy that drives it. PA-DCT maintains discounted dual Gaussian posteriors over both quality and cost per (provider, task-type) bucket, treating cost as a learned signal rather than a known scalar so that price shocks are detected the same way quality shocks are. We construct \emph{402Pilot-Bench}, a reproducible replay benchmark of $824$ tasks across coding, multi-hop QA, and web search; five heterogeneous LLM providers including adversarial and flaky variants at the same price tier; three calibrated market scenarios (stationary, mid-tier outage, premium price promotion); and $30$ seeds $\times$ $10{,}000$ rounds per cell, totaling $20{,}575$ pre-generated frozen LLM responses. PA-DCT outperforms five baselines under both non-stationary scenarios with paired-bootstrap significance at $p<10^{-3}$, while paying a small honest exploration cost in the stationary baseline. A four-component ablation isolates which classical bandit ingredients transfer to the agent-native regime, and shows that no single metric reveals all four; we recommend a multi-metric evaluation framework for future work in this setting.
